\documentclass[11pt]{article}

\usepackage{amsmath,amssymb,amsthm}
\usepackage{geometry}
\geometry{a4paper,margin=1in}

\newtheorem{theorem}{Theorem}
\newtheorem{proposition}[theorem]{Proposition}
\newtheorem{lemma}[theorem]{Lemma}
\newtheorem{remark}[theorem]{Remark}

\title{Verified Notes on the Dual Weight Distribution of the Code in Theorem 4}
\author{ShuYu}
\date{}

\begin{document}

\maketitle

\begin{abstract}
Starting from the source Theorem~4 and Table~\ref{tab:source-table1}, we revisit the dual code of the binary minimal linear code $C_{f,g}$. We first reconcile the parameters and note that the table yields a code with seven nonzero weights. Using the classical MacWilliams identity as input, we then give the exact MacWilliams--Krawtchouk formula for the full dual weight distribution, prove that $A_1^{\perp}=A_2^{\perp}=0$, derive the explicit formulas for $A_3^{\perp}$, $A_4^{\perp}$, and $A_5^{\perp}$, prove that $A_3^{\perp}>0$ under the admissible parameter range, and hence obtain $d(C_{f,g}^{\perp})=3$. Finally, we show that if $s$ is even, then the dual weight distribution is symmetric. All new code-specific results proved in this note are formally verified in Lean~4. Only the general MacWilliams identity, used as a standard cited theorem, and the final qualitative remark about $C^{\perp}$ not being a few-weight code remain outside the current Lean formalization.
\end{abstract}

\section{Source Statement}

We begin by recording the source statement visible in the supplied image. The definitions of $f$, $g$, and the code $C_{f,g}$ come from the preceding Theorem~3 in the source paper.

\begin{theorem}[Source Theorem 4]
Let $s=y$ and $z=1$. If $m\ge 6$ and $2 \le s \le 2^{t-1}-1$, then the code $C_{f,g}$ in Theorem~3 is a $[2^m-1,m+2]$ binary minimal linear code, and its weight distribution is given by Table~\ref{tab:source-table1}. Furthermore, if $s \le 2^{t-2}$, then
\begin{equation}
\frac{w_{\min}}{w_{\max}} \le \frac{1}{2}.
\end{equation}
\end{theorem}

\section{Parameters from Source Theorem 4 and Table 1}

Let $q=2^t$. Table~\ref{tab:source-table1} gives the weights and multiplicities of the code $C=C_{f,g}$:
\begin{table}[h]
\centering
\caption{the weight distribution of $C_{f,g}$ in Theorem~4.}
\label{tab:source-table1}
\begin{tabular}{c|c}
$w$ & $A_w$ \\
\hline
$0$ & $1$ \\
$s(q-1)$ & $2$ \\
$(2s-2)(q-1)$ & $1$ \\
$2^{m-1}$ & $2^m-1$ \\
$2^{m-1}-s$ & $2(q+1-s)(q-1)$ \\
$2^{m-1}-(2s-2)$ & $(q+3-2s)(q-1)$ \\
$2^{m-1}+q-s$ & $2s(q-1)$ \\
$2^{m-1}+q-(2s-2)$ & $(2s-2)(q-1)$
\end{tabular}
\end{table}

Thus there are seven nonzero weights, or eight weights in total when the zero weight is included.

The total number of codewords obtained from the table is
\begin{equation}
\sum_w A_w
=1+2+1+(2^m-1)+2(q+1-s)(q-1)+(q+3-2s)(q-1)+2s(q-1)+(2s-2)(q-1)
=2^m+3q^2.
\end{equation}

Since $C$ is a binary linear code, $|C|$ must be a power of $2$. Therefore the integer
\begin{equation}
2^m+3q^2=2^m+3\cdot 2^{2t}
\end{equation}
must be a power of $2$.

\begin{proposition}
The parameters satisfy
\begin{equation}
m=2t.
\end{equation}
Consequently,
\begin{equation}
|C|=2^m+3q^2=4q^2=2^{m+2}.
\end{equation}
\end{proposition}

\begin{proof}
If $m<2t$, then
\begin{equation}
|C|=2^m\bigl(1+3\cdot 2^{2t-m}\bigr),
\end{equation}
and the factor $1+3\cdot 2^{2t-m}$ is an odd integer strictly larger than $1$, so $|C|$ cannot be a power of $2$.

If $m>2t$, then
\begin{equation}
|C|=2^{2t}\bigl(2^{m-2t}+3\bigr),
\end{equation}
and the factor $2^{m-2t}+3$ is again an odd integer strictly larger than $1$, so $|C|$ cannot be a power of $2$.

Hence neither $m<2t$ nor $m>2t$ is possible, and therefore $m=2t$. Substituting this into $|C|=2^m+3q^2$ gives $|C|=4q^2=2^{m+2}$.
\end{proof}

This parameter reconciliation is formally verified in Lean~4.

After this derivation, the weights may be rewritten as
\begin{equation}
2^{m-1}=q^2/2,
\end{equation}
so the nonzero weights become
\begin{equation}
s(q-1),\ (2s-2)(q-1),\ q^2/2,\ q^2/2-s,\ q^2/2-(2s-2),\ q^2/2+q-s,\ q^2/2+q-(2s-2),
\end{equation}
and the length is
\begin{equation}
n=2^m-1=q^2-1.
\end{equation}

\section{The dual weight enumerator}

Let
\begin{equation}
W_C(x)=\sum_{i=0}^{n}A_i x^i
\end{equation}
be the weight enumerator of $C$, and let
\begin{equation}
W_{C^{\perp}}(x)=\sum_{j=0}^{n}A_j^{\perp}x^j
\end{equation}
be the weight enumerator of the dual code.

\begin{theorem}
For $0\le j\le n=q^2-1$,
\begin{equation}
A_j^{\perp}=\frac1{4q^2}\sum_{w} A_w K_j(w),
\end{equation}
where the sum runs over the eight weights listed above and
\begin{equation}
K_j(w)=\sum_{\ell=0}^{j}(-1)^{\ell}\binom{w}{\ell}\binom{n-w}{j-\ell}
\end{equation}
is the binary Krawtchouk polynomial.

Equivalently,
\begin{equation}
W_{C^{\perp}}(x)=\frac{(1+x)^n}{4q^2}W_C\!\left(\frac{1-x}{1+x}\right).
\end{equation}
\end{theorem}

\begin{proof}
Since $|C|=4q^2$, the binary MacWilliams identity gives
\begin{equation}
W_{C^{\perp}}(x)=\frac{(1+x)^n}{|C|}W_C\!\left(\frac{1-x}{1+x}\right)
=\frac{(1+x)^n}{4q^2}W_C\!\left(\frac{1-x}{1+x}\right).
\end{equation}
Extracting the coefficient of $x^j$ yields the Krawtchouk expansion.
\end{proof}

\begin{remark}
The theorem above is used as a classical coding-theoretic input and is cited rather than formalized here. What is formally verified in Lean~4 are the new code-specific consequences extracted from this formula in the rest of the note.
\end{remark}

\section{The first dual coefficients}

Let
\begin{equation}
M_r=\sum_w A_w w^r
\end{equation}
denote the $r$-th power moment of $C$.

\begin{lemma}
The first two power moments are
\begin{equation}
M_1=2q^2(q^2-1),
\qquad
M_2=q^4(q^2-1).
\end{equation}
\end{lemma}

\begin{proof}
This follows by direct substitution of the eight entries from Table~\ref{tab:source-table1} and straightforward simplification.
\end{proof}

These two power-moment identities are formally verified in Lean~4.

\begin{proposition}
The dual code satisfies
\begin{equation}
A_1^{\perp}=0,
\qquad
A_2^{\perp}=0.
\end{equation}
\end{proposition}

\begin{proof}
Since $|C|=4q^2=2^{m+2}$, the dimension is $k=m+2$, and $n=q^2-1$.

By the first two Pless power moments,
\begin{equation}
M_1=2^{k-1}(n-A_1^{\perp}),
\end{equation}
and
\begin{equation}
M_2=2^{k-2}\bigl(n(n+1)-2nA_1^{\perp}+2A_2^{\perp}\bigr).
\end{equation}
Using $2^k=4q^2$, $n=q^2-1$, $M_1=2q^2(q^2-1)$, and $M_2=q^4(q^2-1)$, the first identity yields $A_1^{\perp}=0$, and then the second yields $A_2^{\perp}=0$.
\end{proof}

The vanishing statements $A_1^{\perp}=0$ and $A_2^{\perp}=0$ are formally verified in Lean~4.

\section{Explicit low-weight multiplicities}

\begin{theorem}
The number of codewords of weight $3$ in $C^{\perp}$ is
\begin{equation}
\begin{aligned}
A_3^{\perp}
=\frac{1}{6}\Bigl(&q^4+(3-6s)q^3+(18s^2-26s+10)q^2 \\
&+(-20s^3+48s^2-40s+12)q+(20s^3-66s^2+72s-26)\Bigr).
\end{aligned}
\end{equation}
\end{theorem}

\begin{proof}
From the third Pless power moment and the already proved equalities $A_1^{\perp}=A_2^{\perp}=0$, one obtains
\begin{equation}
M_3=2^{k-3}\bigl(n^2(n+3)-6A_3^{\perp}\bigr).
\end{equation}
Substituting the Table~\ref{tab:source-table1} weights and multiplicities into $M_3$ and simplifying gives the stated formula.
\end{proof}

This formula is formally verified in Lean~4. In the formal development, it is also checked independently from the Krawtchouk-polynomial expression.

\begin{proposition}
For the admissible range $q=2^t$, $t\ge 3$, and $2\le s\le q/2-1$, one has
\begin{equation}
A_3^{\perp}>0.
\end{equation}
Hence
\begin{equation}
d(C^{\perp})=3.
\end{equation}
\end{proposition}

\begin{proof}
Set
\begin{equation}
r=q-2s.
\end{equation}
Then $2\le r\le q-4$. Rewriting the formula for $A_3^{\perp}$ gives
\begin{equation}
A_3^{\perp}=\frac{q-1}{12}F(q,r),
\end{equation}
where
\begin{equation}
F(q,r)=3q^2r-6qr^2+5r^3+9q^2-34qr+33r^2-44q+72r+52.
\end{equation}
For fixed $r$, we have
\begin{equation}
\frac{\partial F}{\partial q}=6qr-6r^2+18q-34r-44.
\end{equation}
Since $q\ge r+4$, this implies
\begin{equation}
\frac{\partial F}{\partial q}
\ge 6r(r+4)-6r^2+18(r+4)-34r-44
=8r+28>0.
\end{equation}
Thus $F(q,r)$ is increasing in $q$, and so
\begin{equation}
F(q,r)\ge F(r+4,r)=2(r^3+4r^2+6r+10)>0.
\end{equation}
Therefore $A_3^{\perp}>0$. Together with $A_1^{\perp}=A_2^{\perp}=0$, this proves $d(C^{\perp})=3$.
\end{proof}

The positivity statement $A_3^{\perp}>0$ and the conclusion $d(C^{\perp})=3$ are formally verified in Lean~4.

\begin{proposition}
The next two coefficients are
\begin{equation}
\begin{aligned}
A_4^{\perp}=\frac1{24}\Bigl(&q^6-8q^5s+36q^4s^2-80q^3s^3+72q^2s^4+4q^5-40q^4s+120q^3s^2 \\
&-16q^2s^3-288qs^4+13q^4-56q^3s-324q^2s^2+1104qs^3+216s^4 \\
&-4q^3+488q^2s-1632qs^2-1008s^3-198q^2+1056qs+1800s^2 \\
&-240q-1440s+424\Bigr),
\end{aligned}
\end{equation}
and
\begin{equation}
\begin{aligned}
A_5^{\perp}=\frac1{120}\Bigl(&q^8-10q^7s+60q^6s^2-200q^5s^3+360q^4s^4-272q^3s^5+5q^7-70q^6s \\
&+360q^5s^2-680q^4s^3-160q^3s^4+1360q^2s^5+20q^6-220q^5s+180q^4s^2 \\
&+2560q^3s^3-4600q^2s^4-2720qs^5+30q^5+508q^4s-4520q^3s^2+4400q^2s^3 \\
&+14240qs^4+1632s^5-264q^4+2672q^3s+2560q^2s^2-30720qs^3-9840s^4 \\
&-440q^3-6080q^2s+32800qs^2+24640s^3+2440q^2-16960qs-31440s^2 \\
&+3360q+20160s-5152\Bigr).
\end{aligned}
\end{equation}
\end{proposition}

\begin{proof}
These formulas follow by direct evaluation of the MacWilliams--Krawtchouk expression for $A_j^{\perp}$ with $j=4,5$. Specifically, using the Krawtchouk polynomials in the variable $D=n-2w$,
\begin{equation}
24 K_4(w) = D^4-(6n-8)D^2+3n(n-2),
\end{equation}
\begin{equation}
120 K_5(w) = D^5-(10n-20)D^3+(15n^2-50n+24)D,
\end{equation}
we substitute the $D$-values for each table weight and verify that
\begin{equation}
\sum_w A_w D_w^4 = 4q^2\bigl(\text{numerator of } 24A_4^{\perp}+n(3n-2)\bigr),
\end{equation}
\begin{equation}
\sum_w A_w D_w^5 = 4q^2\bigl(\text{numerator of } 120A_5^{\perp}+(10q^2-30)\cdot\text{numerator of } 6A_3^{\perp}\bigr).
\end{equation}
Both identities are polynomial equalities in $q$ and $s$, verified formally in Lean by the \texttt{ring} tactic.
\end{proof}

\section{Symmetry for even $s$}

\begin{proposition}
If $s$ is even, then
\begin{equation}
A_j^{\perp}=A_{n-j}^{\perp}
\qquad
(0\le j\le n).
\end{equation}
\end{proposition}

\begin{proof}
When $s$ is even, every weight appearing in Table~\ref{tab:source-table1} is even: $q^2/2$ is even because $q=2^t$ with $t\ge 3$, and the remaining seven weights are visibly even as well. Hence every codeword of $C$ has even Hamming weight, so $C$ is an even code. Therefore the all-ones vector $\mathbf{1}$ lies in $C^{\perp}$. The map
\begin{equation}
c\longmapsto c+\mathbf{1}
\end{equation}
is then a bijection on $C^{\perp}$ sending words of weight $j$ to words of weight $n-j$, and the symmetry follows.
\end{proof}

This symmetry statement is formally verified in Lean~4.

\begin{remark}
The MacWilliams formula gives the complete dual weight distribution. In general, $C^{\perp}$ is not a few-weight code, so the coefficients $A_j^{\perp}$ are most naturally described by the Krawtchouk expansion rather than by a short finite list. This is a qualitative observation rather than a formalized theorem in the current Lean development.
\end{remark}

\end{document}
